% Ad Atticum 13.25 --- headnote
% ad-atticum-13-25, 12 July 45 BC, Tusculum

\textit{Cicero to Atticus, written at the Tusculan villa on 12 July
45 BC --- Perseus dateline \textit{Scr. in Tusculano iv Id. Quint.
a. 709 (45)}. Two sections of brisk practical business: first,
the timing of Cicero's coming to Rome and the choreography of
not making it look as if he were going up on the Ides only to
``escort'' Brutus to the city or to be present for Brutus's
will-signing; second, the still-unresolved question of whether
Atticus is willing to take the risk of presenting the
\textit{Academica} to Varro as their dedicatee. (The
manuscript transmission preserves no section 2 of this letter,
so the Latin --- and hence this translation --- runs
\textsc{1}, \textsc{3}.)}

\textit{The texture is intimate, telegraphic, half-anxious. Four
short Greek phrases punctuate it: \textit{asaphesteros}
(``rather too obscure,'' a writer's self-criticism that Atticus's
already-read letter exposed); \textit{deinos an\=er}, ``a
formidable man,'' followed by the half-line of Homer
(\textit{Il.} 11.654) \textit{tacha ken kai anaition aiti\=o\=oto}
--- ``he might well bring a charge even against one in no way to
blame'' --- which is Cicero's image of Varro the prickly
dedicatee, ready to take offense whatever the book does; and
\textit{periochas}, the ``outlines'' or running summaries that
Tiro can take down at speed, contrasted with the syllable-by-
syllable dictation Cicero just gave to Spintharus for the
labour-intensive Varro letter. The closing flourish ---
\textit{o Academiam volaticam et sui similem!}, ``O the flighty
Academy, true to itself!'' --- catches the joke whole: the New
Academy's signature is suspension of judgement, and Cicero is
suspending judgement on which of two Academics, Varro or Brutus,
to give the book to. The crux \textit{ad tabulam} I render
``for the auction'' (Shackleton Bailey's reading; \textit{tabula}
as the auctioneer's notice-board), though some take it of an
accounts-tablet for the will.}
